\documentclass[10pt]{article}
\usepackage{amsmath,amsfonts,amssymb,mathtools,mathrsfs,bm,bbm,array,xcolor,geometry}
\allowdisplaybreaks
\setlength{\parindent}{0pt}
\setlength{\parskip}{4pt}
\emergencystretch=4em

\geometry{paperwidth=215mm,paperheight=290mm,hmargin=15mm,vmargin=16mm}
\begin{document}
\section*{P0310 Allocating Relief Supplies Using Regional Urgency}
\subsection*{Mathematical formulation}
\begin{align}
\gamma_g(t)&=\frac{\Omega_g(t)}{\sum_{h\in G}\Omega_h(t)},\\
F_g^1(X)&=\frac{\sum_{l\in L,m\in M,i\in I_g}X_{mi}^l(t)}{D_g},\\
F_g^2(X)&=\sum_{l\in L,m\in M,i\in I_g}
  \bigl[CS_m^l u_i^1(t)+CT_{mi}^l\bigr]X_{mi}^l(t),\\
F_g(X)&=\omega_r\frac{F_g^1(X)-\underline F_g^1}{\overline F_g^1-\underline F_g^1}
-\omega_c\frac{F_g^2(X)-\underline F_g^2}{\overline F_g^2-\underline F_g^2}.
\end{align}

\begin{align}
\max_X\quad &\sum_{g\in G}\gamma_g(t)F_g(X)\tag{20}\\
\text{s.t.}\quad
&\sum_{m\in M}X_{mi}^l(t)\le D_i^l(t)
 &&g\in G,\ i\in I_g,\ l\in L,\tag{22}\\
&\sum_{g\in G,i\in I_g}X_{mi}^l(t)\le Q_m^l(t-1|t-1)
 &&m\in M,\ l\in L,\tag{23}\\
&\sum_{g\in G,i\in I_g,l\in L}UNV^lX_{mi}^l(t)
\le\sum_{v\in V}UNL_vNV_v^m(t) &&m\in M,\tag{24}\\
&X_{mi}^l(t)\ge0 &&m\in M,\ g\in G,\ i\in I_g,\ l\in L.\tag{25}
\end{align}

\subsection*{Definitions and model scope}
For one fixed decision period, a relief coordinator allocates several supply types $L$ from given distribution centres $M$ to affected regions. Given regional groups $G$ contain disjoint region memberships $I_g$. Demand forecasts, regional grouping and urgency assessments are already available; they are not decisions in this allocation. Each region and supply type has a nonnegative demand $D_i^l(t)$, and each centre has nonnegative beginning stock $Q_m^l(t-1|t-1)$ for each supply type. The coordinator chooses continuous nonnegative shipment quantities $X_{mi}^l(t)$ from every centre to every region for every supply type.

For each region and supply type, add shipments from all centres; the total cannot exceed that region's demand. For each centre and supply type, add shipments to all regional groups and their members; the total cannot exceed that centre's beginning stock. Thus supplies cannot exceed the recorded local demand or be taken from stock that is not available. A group with zero total demand receives no shipment and is not assigned a demand-fulfilment ratio requiring division by zero.

Each supply type has a given volume per unit $UNV^l$. At each centre, multiply every shipment by its supply's unit volume and add over all supplies and destination regions. The total cannot exceed the sum, over the given vehicle types $V$, of each type's vehicle capacity $UNL_v$ times its available trips at that centre $NV_v^m(t)$. These trip counts are supplied for the current period; the coordinator does not separately build vehicle routes or decide which vehicle carries an individual shipment.

For each group with positive total demand $D_g$, its fulfilment measure $F_g^1(X)$ is the sum of shipments of all supply types from all centres to all regions in the group, divided by the sum of that group's demands over supplies and regions. Its delivery cost $F_g^2(X)$ is calculated shipment by shipment: multiply the given unit setup-cost coefficient $CS_m^l$ by the region's given time since the last arrival of supplies $u_i^1(t)$, add the given unit transport cost $CT_{mi}^l$, and multiply this whole unit amount by the shipment quantity. Add across the group's supplies, centres and regions; the setup-cost coefficient here is not a fixed charge merely for opening a route.

For each group and for each of its two measures, given lower reference levels $\underline F_g^j$ and upper reference levels $\overline F_g^j$ have the upper level strictly above the lower. Normalize each measure by subtracting its lower reference level and dividing by the difference between its upper and lower levels. The group's combined score $F_g(X)$ is the normalized fulfilment multiplied by the nonnegative fulfilment weight $\omega_r$, minus normalized cost multiplied by the nonnegative cost weight $\omega_c$. These two weights sum to one; reference levels normalize the score and do not separately impose new shipment bounds.

The given group urgency assessment $\Omega_g(t)$ is divided by the sum of the assessments over the given groups to obtain its urgency weight $\gamma_g(t)$; that denominator must be nonzero. The coordinator seeks the greatest sum of urgency weight times combined score over the participating groups. All forecasts, urgency assessments, normalization limits, costs, stocks and available transport capacities remain fixed during this choice. Later supplier replenishment or inventory updates may use the resulting shipments, but their separate aims do not alter this single-period allocation.

\subsection*{Inline notation and relations}
$t$; $L$; $M$; $G$; $I_g$; $g$; $V$; $D_i^l(t)\ge0$; $Q_m^l(t-1|t-1)\ge0$; $UNV^l$; $UNL_v$; $NV_v^m(t)$; $CS_m^l$; $CT_{mi}^l$; $u_i^1(t)$; $\Omega_g(t)$; $X_g(t)$; $\overline F_g^j>\underline F_g^j$; $j=1,2$; $\omega_r,\omega_c\ge0$; $\omega_r+\omega_c=1$; $D_g=\sum_{l\in L,i\in I_g}D_i^l(t)>0$; $\sum_g\Omega_g(t)\ne0$; $X_{mi}^l(t)\ge0$; $m$; $i$; $l$; $v$.
\end{document}